% Options for packages loaded elsewhere
% Options for packages loaded elsewhere
\PassOptionsToPackage{unicode}{hyperref}
\PassOptionsToPackage{hyphens}{url}
%
\documentclass[
  letterpaper,
  oneside,
  open=any]{scrbook}
\usepackage{xcolor}
\usepackage{amsmath,amssymb}
\setcounter{secnumdepth}{5}
\usepackage{iftex}
\ifPDFTeX
  \usepackage[T1]{fontenc}
  \usepackage[utf8]{inputenc}
  \usepackage{textcomp} % provide euro and other symbols
\else % if luatex or xetex
  \usepackage{unicode-math} % this also loads fontspec
  \defaultfontfeatures{Scale=MatchLowercase}
  \defaultfontfeatures[\rmfamily]{Ligatures=TeX,Scale=1}
\fi
\usepackage{lmodern}
\ifPDFTeX\else
  % xetex/luatex font selection
\fi
% Use upquote if available, for straight quotes in verbatim environments
\IfFileExists{upquote.sty}{\usepackage{upquote}}{}
\IfFileExists{microtype.sty}{% use microtype if available
  \usepackage[]{microtype}
  \UseMicrotypeSet[protrusion]{basicmath} % disable protrusion for tt fonts
}{}
\makeatletter
\@ifundefined{KOMAClassName}{% if non-KOMA class
  \IfFileExists{parskip.sty}{%
    \usepackage{parskip}
  }{% else
    \setlength{\parindent}{0pt}
    \setlength{\parskip}{6pt plus 2pt minus 1pt}}
}{% if KOMA class
  \KOMAoptions{parskip=half}}
\makeatother
% Make \paragraph and \subparagraph free-standing
\makeatletter
\ifx\paragraph\undefined\else
  \let\oldparagraph\paragraph
  \renewcommand{\paragraph}{
    \@ifstar
      \xxxParagraphStar
      \xxxParagraphNoStar
  }
  \newcommand{\xxxParagraphStar}[1]{\oldparagraph*{#1}\mbox{}}
  \newcommand{\xxxParagraphNoStar}[1]{\oldparagraph{#1}\mbox{}}
\fi
\ifx\subparagraph\undefined\else
  \let\oldsubparagraph\subparagraph
  \renewcommand{\subparagraph}{
    \@ifstar
      \xxxSubParagraphStar
      \xxxSubParagraphNoStar
  }
  \newcommand{\xxxSubParagraphStar}[1]{\oldsubparagraph*{#1}\mbox{}}
  \newcommand{\xxxSubParagraphNoStar}[1]{\oldsubparagraph{#1}\mbox{}}
\fi
\makeatother

\usepackage{color}
\usepackage{fancyvrb}
\newcommand{\VerbBar}{|}
\newcommand{\VERB}{\Verb[commandchars=\\\{\}]}
\DefineVerbatimEnvironment{Highlighting}{Verbatim}{commandchars=\\\{\}}
% Add ',fontsize=\small' for more characters per line
\usepackage{framed}
\definecolor{shadecolor}{RGB}{241,243,245}
\newenvironment{Shaded}{\begin{snugshade}}{\end{snugshade}}
\newcommand{\AlertTok}[1]{\textcolor[rgb]{0.68,0.00,0.00}{#1}}
\newcommand{\AnnotationTok}[1]{\textcolor[rgb]{0.37,0.37,0.37}{#1}}
\newcommand{\AttributeTok}[1]{\textcolor[rgb]{0.40,0.45,0.13}{#1}}
\newcommand{\BaseNTok}[1]{\textcolor[rgb]{0.68,0.00,0.00}{#1}}
\newcommand{\BuiltInTok}[1]{\textcolor[rgb]{0.00,0.23,0.31}{#1}}
\newcommand{\CharTok}[1]{\textcolor[rgb]{0.13,0.47,0.30}{#1}}
\newcommand{\CommentTok}[1]{\textcolor[rgb]{0.37,0.37,0.37}{#1}}
\newcommand{\CommentVarTok}[1]{\textcolor[rgb]{0.37,0.37,0.37}{\textit{#1}}}
\newcommand{\ConstantTok}[1]{\textcolor[rgb]{0.56,0.35,0.01}{#1}}
\newcommand{\ControlFlowTok}[1]{\textcolor[rgb]{0.00,0.23,0.31}{\textbf{#1}}}
\newcommand{\DataTypeTok}[1]{\textcolor[rgb]{0.68,0.00,0.00}{#1}}
\newcommand{\DecValTok}[1]{\textcolor[rgb]{0.68,0.00,0.00}{#1}}
\newcommand{\DocumentationTok}[1]{\textcolor[rgb]{0.37,0.37,0.37}{\textit{#1}}}
\newcommand{\ErrorTok}[1]{\textcolor[rgb]{0.68,0.00,0.00}{#1}}
\newcommand{\ExtensionTok}[1]{\textcolor[rgb]{0.00,0.23,0.31}{#1}}
\newcommand{\FloatTok}[1]{\textcolor[rgb]{0.68,0.00,0.00}{#1}}
\newcommand{\FunctionTok}[1]{\textcolor[rgb]{0.28,0.35,0.67}{#1}}
\newcommand{\ImportTok}[1]{\textcolor[rgb]{0.00,0.46,0.62}{#1}}
\newcommand{\InformationTok}[1]{\textcolor[rgb]{0.37,0.37,0.37}{#1}}
\newcommand{\KeywordTok}[1]{\textcolor[rgb]{0.00,0.23,0.31}{\textbf{#1}}}
\newcommand{\NormalTok}[1]{\textcolor[rgb]{0.00,0.23,0.31}{#1}}
\newcommand{\OperatorTok}[1]{\textcolor[rgb]{0.37,0.37,0.37}{#1}}
\newcommand{\OtherTok}[1]{\textcolor[rgb]{0.00,0.23,0.31}{#1}}
\newcommand{\PreprocessorTok}[1]{\textcolor[rgb]{0.68,0.00,0.00}{#1}}
\newcommand{\RegionMarkerTok}[1]{\textcolor[rgb]{0.00,0.23,0.31}{#1}}
\newcommand{\SpecialCharTok}[1]{\textcolor[rgb]{0.37,0.37,0.37}{#1}}
\newcommand{\SpecialStringTok}[1]{\textcolor[rgb]{0.13,0.47,0.30}{#1}}
\newcommand{\StringTok}[1]{\textcolor[rgb]{0.13,0.47,0.30}{#1}}
\newcommand{\VariableTok}[1]{\textcolor[rgb]{0.07,0.07,0.07}{#1}}
\newcommand{\VerbatimStringTok}[1]{\textcolor[rgb]{0.13,0.47,0.30}{#1}}
\newcommand{\WarningTok}[1]{\textcolor[rgb]{0.37,0.37,0.37}{\textit{#1}}}

\usepackage{longtable,booktabs,array}
\usepackage{calc} % for calculating minipage widths
% Correct order of tables after \paragraph or \subparagraph
\usepackage{etoolbox}
\makeatletter
\patchcmd\longtable{\par}{\if@noskipsec\mbox{}\fi\par}{}{}
\makeatother
% Allow footnotes in longtable head/foot
\IfFileExists{footnotehyper.sty}{\usepackage{footnotehyper}}{\usepackage{footnote}}
\makesavenoteenv{longtable}
\usepackage{graphicx}
\makeatletter
\newsavebox\pandoc@box
\newcommand*\pandocbounded[1]{% scales image to fit in text height/width
  \sbox\pandoc@box{#1}%
  \Gscale@div\@tempa{\textheight}{\dimexpr\ht\pandoc@box+\dp\pandoc@box\relax}%
  \Gscale@div\@tempb{\linewidth}{\wd\pandoc@box}%
  \ifdim\@tempb\p@<\@tempa\p@\let\@tempa\@tempb\fi% select the smaller of both
  \ifdim\@tempa\p@<\p@\scalebox{\@tempa}{\usebox\pandoc@box}%
  \else\usebox{\pandoc@box}%
  \fi%
}
% Set default figure placement to htbp
\def\fps@figure{htbp}
\makeatother





\setlength{\emergencystretch}{3em} % prevent overfull lines

\providecommand{\tightlist}{%
  \setlength{\itemsep}{0pt}\setlength{\parskip}{0pt}}



 


\usepackage{booktabs}
\usepackage{caption}
\usepackage{longtable}
\usepackage{colortbl}
\usepackage{array}
\usepackage{anyfontsize}
\usepackage{multirow}
\makeatletter
\@ifpackageloaded{bookmark}{}{\usepackage{bookmark}}
\makeatother
\makeatletter
\@ifpackageloaded{caption}{}{\usepackage{caption}}
\AtBeginDocument{%
\ifdefined\contentsname
  \renewcommand*\contentsname{Table of contents}
\else
  \newcommand\contentsname{Table of contents}
\fi
\ifdefined\listfigurename
  \renewcommand*\listfigurename{List of Figures}
\else
  \newcommand\listfigurename{List of Figures}
\fi
\ifdefined\listtablename
  \renewcommand*\listtablename{List of Tables}
\else
  \newcommand\listtablename{List of Tables}
\fi
\ifdefined\figurename
  \renewcommand*\figurename{Figure}
\else
  \newcommand\figurename{Figure}
\fi
\ifdefined\tablename
  \renewcommand*\tablename{Table}
\else
  \newcommand\tablename{Table}
\fi
}
\@ifpackageloaded{float}{}{\usepackage{float}}
\floatstyle{ruled}
\@ifundefined{c@chapter}{\newfloat{codelisting}{h}{lop}}{\newfloat{codelisting}{h}{lop}[chapter]}
\floatname{codelisting}{Listing}
\newcommand*\listoflistings{\listof{codelisting}{List of Listings}}
\makeatother
\makeatletter
\makeatother
\makeatletter
\@ifpackageloaded{caption}{}{\usepackage{caption}}
\@ifpackageloaded{subcaption}{}{\usepackage{subcaption}}
\makeatother
\usepackage{bookmark}
\IfFileExists{xurl.sty}{\usepackage{xurl}}{} % add URL line breaks if available
\urlstyle{same}
\hypersetup{
  pdftitle={Phenotyping in the OMOP Common Data Model},
  pdfauthor={Martí Català; Anna Saura-Lazaro; Edward Burn},
  hidelinks,
  pdfcreator={LaTeX via pandoc}}


\title{Phenotyping in the OMOP Common Data Model}
\author{Martí Català \and Anna Saura-Lazaro \and Edward Burn}
\date{2026-06-03}
\begin{document}
\frontmatter
\maketitle

\renewcommand*\contentsname{Table of contents}
{
\setcounter{tocdepth}{2}
\tableofcontents
}

\mainmatter
\bookmarksetup{startatroot}

\chapter*{Preface}\label{preface}
\addcontentsline{toc}{chapter}{Preface}

\markboth{Preface}{Preface}

\section*{Is this book for me?}\label{is-this-book-for-me}
\addcontentsline{toc}{section}{Is this book for me?}

\markright{Is this book for me?}

This book is part of the collection of \textbf{Tidy R programming in
OMOP} see accompanying books:

\begin{enumerate}
\def\labelenumi{(\arabic{enumi})}
\tightlist
\item
  \href{https://oxford-pharmacoepi.github.io/Tidy-R-programming-with-OMOP/}{Tidy
  R programming with the OMOP common data model}
\item
  \href{https://oxford-pharmacoepi.github.io/Tidy-R-packages-with-OMOP/}{Tidy
  R packages with the OMOP common data model}
\item
  \href{https://oxford-pharmacoepi.github.io/DrugUtilisation-in-OMOP-CDM/}{Drug
  Utilisation Studies in the OMOP Common Data Model}
\item
  \href{https://oxford-pharmacoepi.github.io/Phenotyping-in-OMOP/}{Phenotyping
  in the OMOP Common Data Model}
\end{enumerate}

We've written this book for anyone interested to conduct studies using
the OMOP Common Data Model and working with databases using a tidyverse
style approach. That is, human centered, consistent, composable, and
inclusive (see \url{https://design.tidyverse.org/unifying.html} for more
details on these principles).

New to R? We recommend you compliment the book with
\href{https://r4ds.had.co.nz/}{R for data science}

New to databases? We recommend you take a look at some web tutorials on
SQL, such as \href{https://sqlbolt.com/}{SQLBolt} or
\href{https://www.sqlzoo.net/wiki/SQL_Tutorial}{SQLZoo}

New to the OMOP CDM? We'd recommend you pare this book with
\href{https://ohdsi.github.io/TheBookOfOhdsi/}{The Book of OHDSI}

New to the Tidy R in OMOP ecosystem? Probably you should start with the
book
\href{https://oxford-pharmacoepi.github.io/Tidy-R-programming-with-OMOP/}{Tidy
R programming with the OMOP common data model} before reading this one.

\section*{How is the book organised?}\label{how-is-the-book-organised}
\addcontentsline{toc}{section}{How is the book organised?}

\markright{How is the book organised?}

TO ADD

\section*{Citation}\label{citation}
\addcontentsline{toc}{section}{Citation}

\markright{Citation}

TO ADD

\section*{License}\label{license}
\addcontentsline{toc}{section}{License}

\markright{License}

\section*{Code}\label{code}
\addcontentsline{toc}{section}{Code}

\markright{Code}

The source code for the book can be found at this
\href{https://github.com/oxford-pharmacoepi/Phenotyping-in-OMOP}{Github
repository}

\section*{Packages version}\label{packages-version}
\addcontentsline{toc}{section}{Packages version}

\markright{Packages version}

This book is rendered using the following version of packages:

\begin{verbatim}
Finding R package dependencies ... Done!
\end{verbatim}

\begin{table}
\fontsize{12.0pt}{14.0pt}\selectfont
\begin{tabular*}{\linewidth}{@{\extracolsep{\fill}}lrl}
\toprule
Package & Version & Link \\ 
\midrule\addlinespace[2.5pt]
CodelistGenerator & 4.0.2 & \href{https://darwin-eu.github.io/CodelistGenerator/}{🔗} \\ 
PhenotypeR & 0.5.0 & \href{https://ohdsi.github.io/PhenotypeR/}{🔗} \\ 
duckdb & 1.5.0 & \href{https://r.duckdb.org/}{🔗} \\ 
omock & 0.6.2 & \href{https://ohdsi.github.io/omock/}{🔗} \\ 
\bottomrule
\end{tabular*}
\end{table}

Note we only included the packages called explicitly in the book.

\bookmarksetup{startatroot}

\chapter{Proposed structure for the
book}\label{proposed-structure-for-the-book}

\begin{itemize}
\tightlist
\item
  Preface
\item
  Introduction

  \begin{itemize}
  \tightlist
  \item
    What is Phenotyping
  \item
    The phenotyping process
  \item
    Clinical description
  \end{itemize}
\item
  OMOP CDM vocabularies

  \begin{itemize}
  \tightlist
  \item
    Introduction
  \item
    Download vocabularies
  \item
    Standard vs Source vocabulary
  \item
    Vocabularies hierarchy
  \item
    Condition based concepts
  \item
    Drug based concepts
  \item
    Measurement based concepts
  \item
    Observation based concepts
  \item
    Procedures based concepts
  \item
    Devices based concepts
  \end{itemize}
\item
  Codes in the OMOP CDM

  \begin{itemize}
  \tightlist
  \item
    Introduction
  \item
    get candidate codelists
  \item
    Create a condition based codelist (example stroke)
  \item
    Create a \ldots.
  \item
    Create a drug based codelist (example chemotheraphy)
  \item
    Create a drug based codelist (sth that we can stratify by brand,
    route and dose unit)
  \end{itemize}
\item
  Create cohorts

  \begin{itemize}
  \tightlist
  \item
    Base cohorts
  \item
    Demographic cohorts
  \item
    Inclusion criteria
  \item
    Set the end date
  \item
    Combine cohorts
  \item
    Example cohorts (obesity + causal inference centered on a
    complicated end date e.g.~venla study)
  \end{itemize}
\item
  Phenotype validation

  \begin{itemize}
  \tightlist
  \item
    Database Feasibility

    \begin{itemize}
    \tightlist
    \item
      Introduction
    \item
      Metadata
    \item
      The person table
    \item
      The observation period
    \item
      Clinical tables
    \item
      Record trends
    \end{itemize}
  \item
    Codelist diagnostics

    \begin{itemize}
    \tightlist
    \item
      Code use
    \item
      Orphan codes
    \item
      Drug exposure diagnostics
    \item
      Measurement diagnostics
    \end{itemize}
  \item
    Cohort diagnostics

    \begin{itemize}
    \tightlist
    \item
      Counts and attrition
    \item
      Characterisation and large scale characterisation
    \item
      Matched comparison
    \item
      Overlap and timing
    \item
      Survival
    \end{itemize}
  \item
    Population diagnostics

    \begin{itemize}
    \tightlist
    \item
      Incidence
    \item
      Prevalence
    \end{itemize}
  \end{itemize}
\item
  Study workflow and final remarks
\end{itemize}

\begin{Shaded}
\begin{Highlighting}[]
\FunctionTok{library}\NormalTok{(CodelistGenerator)}
\end{Highlighting}
\end{Shaded}

\begin{verbatim}
Warning: package 'CodelistGenerator' was built under R version 4.4.3
\end{verbatim}

\begin{Shaded}
\begin{Highlighting}[]
\FunctionTok{library}\NormalTok{(PhenotypeR)}
\FunctionTok{library}\NormalTok{(omock)}
\FunctionTok{library}\NormalTok{(duckdb)}
\end{Highlighting}
\end{Shaded}

\begin{verbatim}
Warning: package 'duckdb' was built under R version 4.4.3
\end{verbatim}

\begin{verbatim}
Loading required package: DBI
\end{verbatim}

\begin{verbatim}
Warning: package 'DBI' was built under R version 4.4.3
\end{verbatim}

\part{Introduction}

\chapter{Phenotyping process}\label{phenotyping-process}

\part{The OMOP CDM vocabularies}

\chapter{Download vocabularies}\label{download-vocabularies}

In this chapter we explain how to download the OMOP CDM vocabularies and
create an OMOP CDM instance with just the vocabulary tables. We will use
this OMOP CDM instance to explore the vocabularies in the next chapters.

The OMOP CDM vocabularies are distributed thought
\href{athena.ohdsi.org}{Athena} this is a platform where you can explore
the different vocabularies, perform searches and explore the hierarchy
of the vocabularies. If you want to read more about
\href{athena.ohdsi.org}{Athena} you can read the chapter \ldots{} from
\emph{The book of OHDSI}.

\section{Download vocabularies}\label{download-vocabularies-1}

\begin{itemize}
\tightlist
\item
  Create an account in Athena
\item
  Prepare the vocabularies that you are interested.
\item
  Download
\item
  Extract CPT4 \emph{note this step is only needed if }
\item
  create the cdm\_reference
\end{itemize}

\part{Codes in the OMOP CDM}

\chapter{Introduction}\label{introduction-1}

In this chapter we will see the different forms of codelists
(\texttt{\textless{}codelist\textgreater{}},
\texttt{\textless{}codelist\_with\_details\textgreater{}} and
\texttt{\textless{}concept\_set\_expression\textgreater{}})

\bookmarksetup{startatroot}

\chapter*{Final remarks}\label{final-remarks}
\addcontentsline{toc}{chapter}{Final remarks}

\markboth{Final remarks}{Final remarks}

\emph{Drug Utilisation Studies in the OMOP Common Data Model} aims to
\ldots{}

\subsection*{Learning more}\label{learning-more}
\addcontentsline{toc}{subsection}{Learning more}

If you find this book useful then joining the \textbf{Tidy R in OMOP}
OHDSI working group will likely be of interest. Building on many of the
concepts and tools seen in this book, the Tidy R in OMOP OHDSI working
group aims to (1) create and promote a unified set of resources to guide
Tidy R development in OMOP and support the OHDSI community, and (2)
establish an overview of available packages relevant to Tidy R
programming in OMOP (tidyverse style packages). If you are interested in
joining the working group then please email
\href{mailto:catalasabate@ohdsi.org}{Martí Català} or
\href{mailto:kolde@ohdsi.org}{Raivo Kolde}.

\section*{Support us}\label{support-us}
\addcontentsline{toc}{section}{Support us}

\markright{Support us}

We encourage you to support this work by either citing the book in your
papers or documentation, recommending it to your colleagues, or starring
the \href{https://github.com/OHDSI/Tidy-R-programming-with-OMOP}{GitHub
repository}, or simply letting us know how it helped you. Most
importantly, please \textbf{use it} in research that results in patient
benefit.


\backmatter


\end{document}
